%!TEX root = main.tex

\begin{abstract}
Este artículo presenta un análisis comparativo de rendimiento de un kernel de suavizado de imágenes (box filter 3x3) implementado bajo tres paradigmas de programación: escalar (secuencial), multinúcleo (hilos CPU con OpenMP) y heterogéneo (offloading a GPU con OpenMP Target). Utilizando una imagen fractal de 1024x1024 píxeles como caso de prueba y el compilador Intel icpx, se evaluó la eficiencia de cada aproximación mediante la medición de tiempos de ejecución y el análisis de la sobrecarga de sincronización y transferencia de datos. Los resultados demuestran que, para cargas de trabajo moderadas, la versión escalar supera en rendimiento debido al bajo costo de inicialización, mientras que las versiones paralelas enfrentan cuellos de botella significativos derivados de la gestión de hilos y el movimiento de memoria entre el host y el acelerador. Se propone el uso de directivas de optimización de datos para mitigar el impacto de la latencia en arquitecturas heterogéneas.

\textbf{Palabras clave:} Suavizado de imágenes, OpenMP, GPU Offloading, Procesamiento Heterogéneo, Rendimiento, Box Filter.
\end{abstract}